Al utilizar protocolos basados en teoría de grupos, como Diffie-Hellman, ElGamal o firmas de Schnorr, nos basamos en un grupo $G$, un elemento $g$ del grupo y un número natural $q$ que son definidos de antemano por alguna entidad externa. Explique qué son estos elementos $g$ y $q$, cuál es la relación entre ellos. Explique también qué debemos verificar para asegurar que la entidad externa no nos está \emph{engañando} con respecto a esa relación. Justifique por qué dicha verificación asegura que la relación entre $g$ y $q$ es la esperada.

\paragraph{Solución.} $g$ es un elemento del grupo $G$ y $q$ es el orden de $\langle g\rangle$, el subgrupo generado por $g$. Es decir, $q$ es el menor natural tal que $g^q$ es la identidad. Para verificar que $q$ es efectivamente el orden del subgrupo generado por $g$, chequeamos que $q$ sea un número primo y que $g^q$ sea la identidad. Esto asegura que $q$ es el orden del del subgrupo generado por $g$ por el teorema de Lagrange: si existiera un número $x\in\{2,\cdots q\}$ tal que $g^x$ es la identidad, $x$ tendría que dividir a $q$. Pero como $q$ es un número primo, tal $x$ no puede existir.
